\documentclass[10.5pt, a4paper]{ctexart}
\usepackage{../common/ncd-style}

\hypersetup{
  pdftitle={新概念丹麦语 1--5 总目录与语法大纲},
  pdfauthor={Jaehaerys},
  pdfsubject={Ny Koncept Dansk Syllabus}
}

\begin{document}

% Cover / Title block
\begin{titlepage}
  \centering
  \vspace*{2.5cm}
  {\color{NCDRed}\rule{\linewidth}{2pt}}\\[1.2cm]
  {\Huge\bfseries\color{NCDNavy} 《新概念丹麦语 1--5》}\\[0.5cm]
  {\LARGE\itshape\color{NCDTeal} Ny Koncept Dansk (NCD 1--5)}\\[0.6cm]
  {\Large\bfseries\color{NCDGray} 总目录与螺旋式语法大纲}\\[1.2cm]
  {\color{NCDRed}\rule{\linewidth}{2pt}}\\[2cm]

  \begin{tcolorbox}[
    colback=NCDBgBlue,
    colframe=NCDNavy,
    width=0.88\linewidth,
    arc=4pt,
    boxrule=0.8pt,
    left=14pt, right=14pt, top=12pt, bottom=12pt
  ]
    \textbf{\color{NCDNavy}编写理念：}借鉴《新概念英语》（\textit{New Concept English}）“螺旋式递进语法大纲 + 高复现率短文 + 固定句型操练”的自学闭环，结合丹麦官方 \textbf{DU3（Danskuddannelse 3）Modul 1--5 / PD3 $\rightarrow$ Studieprøven} 考核标准与丹麦语核心特性（V2 语序、从句语序、软 d 与吞音、屈折词形、情态小品词），打造专供自主学习的系统化丹麦语教材（共 5 册、104 课）。
  \end{tcolorbox}

  \vfill
  {\large\color{NCDGray} 版本：2026.11 \quad|\quad 全套规划：NCD 1--5（104 课）}
\end{titlepage}

\tableofcontents
\clearpage

% =====================================================================
\section{全套五册定位总览}
% =====================================================================

\renewcommand{\arraystretch}{1.35}
{\small
\begin{tabularx}{\linewidth}{@{} l >{\raggedright\arraybackslash}p{2.9cm} >{\raggedright\arraybackslash}p{2.2cm} >{\raggedright\arraybackslash}p{2.1cm} >{\raggedright\arraybackslash}p{3.0cm} >{\raggedright\arraybackslash}X @{}}
\toprule
\textbf{分册} & \textbf{书名（丹/中）} & \textbf{对标新概念} & \textbf{DU3 / CEFR} & \textbf{体裁与篇幅} & \textbf{核心能力靶点} \\
\midrule
\textbf{NCD 1} & \textit{Første Skridt}\newline （起步篇 · 24 课） & NCE 1\newline \textit{(First Things First)} & \textbf{Modul 1}\newline (A1 $\rightarrow$ A2-) & \textbf{情景对话 + 短独白}\newline （60--100 词） & 发音与吞音感知、名词 \textit{en/et} 与定指复数、现在时/情态动词/高频过去时、基础 V2 语序、疑问词提问、Mindmap 独白 \\
\textbf{NCD 2} & \textit{Fra Sætning til Tekst}\newline （进阶篇 · 24 课） & NCE 2 前半\newline \textit{(Practice \& Progress)} & \textbf{Modul 2}\newline (A2) & \textbf{叙事短文 + 邮件 + 多人对比}\newline （100--180 词） & 从句语序（\textit{ikke/altid} 前置）、现在完成时 vs 过去时切换、核心从句连词、\textit{som/der}、70+ 词邮件写作与阅读四大题型 \\
\textbf{NCD 3} & \textit{Hverdag og Samfund}\newline （社会篇 · 20 课） & NCE 2 后半 $\rightarrow$ NCE 3 前半 & \textbf{Modul 3}\newline (B1) & \textbf{记叙与社会制度说明文}\newline （200--300 词） & \textit{-s} 被动 vs \textit{blive} 被动、介词 + 形式宾语 \textit{det} + \textit{at} 从句、分词定语、省略 \textit{hvis} 条件倒装、复合名词构词法 \\
\textbf{NCD 4} & \textit{Holdning og Debat}\newline （思辨篇 · 20 课） & NCE 3\newline \textit{(Developing Skills)} & \textbf{Modul 4}\newline (B1+ $\rightarrow$ B2-) & \textbf{观点专栏、辩论与读者来信}\newline （300--450 词） & 情态小品词（\textit{jo, da, nok, vist, lige, vel}）、虚拟语气、动词名词化、口语转书面语、\textit{Debatindlæg} 写作 \\
\textbf{NCD 5} & \textit{Tanke og Perspektiv}\newline （精研篇 · 16 课） & NCE 4\newline \textit{(Fluency)} & \textbf{Modul 5 (PD3)} $\rightarrow$ \textbf{Modul 6}\newline (B2 $\rightarrow$ C1) & \textbf{社论、图表评论、思辨长文}\newline （450--650 词） & 统计图表三步评论法、低调反讽（\textit{Underdrivelse}）、学术引述立场动词、多层嵌套长句与多源材料综合论述 \\
\bottomrule
\end{tabularx}
}

% =====================================================================
\section{30 天速通规划（NCD 1 + NCD 2 备考 DU3 Modul 2 专用）}
% =====================================================================

结合\textbf{1 个月内学完 NCD 1 与 NCD 2 并冲刺 DU3 Modul 2}的目标，推荐采用“查漏补缺 + 双轨推进”的 30 天节奏：

\begin{tipbox}{30 天三阶段备考节奏表}
\begin{itemize}
  \item \textbf{第 1--10 天｜NCD 1 速通固基期（每天 2--3 课，共 24 课）：}扫除口腔肌肉与发音盲区（软 d、不发音辅音、吞音连读）；把“\textbf{主句 V2 倒装}”、“\textbf{中心状语 \textit{ikke/altid/også} 位置}”、“\textbf{\textit{en/et} 搭配 \textit{hvilken/hvilket}、\textit{nogen/noget}}”以及“\textbf{20 个高频不规则动词过去时}”练成条件反射。Unit 1--5 每单元第 4 课独白可直接作为口试 \textit{Opgave 1}（Mindmap 陈述）模板，对话问句支撑 \textit{Opgave 2}（信息差提问）。
  \item \textbf{第 11--26 天｜NCD 2 核心攻坚期（每天 1.5 课，共 24 课）：}深度对标 Modul 2 四大核心考点：
    \begin{enumerate}
      \item \textbf{L1--L6}：私人邮件写作（70+ 词）+ 现在完成时 \textit{har/er} 区分 + 完成时与过去时黄金切换 + \textit{fordi/da/når} 从句；
      \item \textbf{L7--L12}：阅读 \textit{Opgave 1 \& 2} 逻辑连贯训练 + \textit{selvom/hvis/mens} + 比较级与过去完成时；
      \item \textbf{L13--L18}：阅读 \textit{Opgave 3 \& 4} 完形与三人物对比 + \textit{som/der} + 篇章副词 \textit{derfor/alligevel} + \textit{sin/sit/sine} 排雷；
      \item \textbf{L19--L24}：长句转述与跨模块衔接，为进入 Modul 3 铺路。
    \end{enumerate}
  \item \textbf{第 27--30 天｜Modul 2 冲刺模考期（4 天）：}每天完成 1 套完整 DU3 Modul 2 计时模拟（阅读 4 题 + 70 词以上邮件写作 + 口语 Mindmap 2 分钟陈述与 6 题信息差提问），回看 NCD 2 中错题对应的语法卡片。
\end{itemize}
\end{tipbox}

% =====================================================================
\section{五册螺旋式语法递进总表（Spiral Grammar Matrix）}
% =====================================================================

与《新概念英语》相同，核心语法不在某几课一次性讲完，而是分五轮螺旋加深：

{\footnotesize
\renewcommand{\arraystretch}{1.3}
\begin{tabularx}{\linewidth}{@{} >{\bfseries\raggedright\arraybackslash}p{1.9cm} *{5}{>{\raggedright\arraybackslash}X} @{}}
\toprule
\textbf{语法维度} & \textbf{NCD 1 (A1--A2-)}\newline Modul 1 & \textbf{NCD 2 (A2)}\newline Modul 2 & \textbf{NCD 3 (B1)}\newline Modul 3 & \textbf{NCD 4 (B1+--B2-)}\newline Modul 4 & \textbf{NCD 5 (B2--C1)}\newline PD3 $\rightarrow$ Modul 6 \\
\midrule
1. 语序与\newline 句法骨架 & 陈述句 SVO；是非/hv- 问句倒装；句首时间/地点触发\textbf{主句 V2 倒装}；中心状语（\textit{ikke/altid}）位置 & \textbf{从句语序（主语 + \textit{ikke/altid} + 动词）}；“从句在前、主句倒装”；间接疑问句（\textit{om / hv-}） & 无 \textit{hvis} 省略条件倒装句；分裂强调句（\textit{Det er X, der...}）；\textbf{介词 + 形式宾语 \textit{det} + \textit{at} 从句} & 口语左移外置与书面紧凑句互转；多层从句嵌套；被动分词前置长定语拆解 & 政论修辞长句、高阶连词（\textit{idet, eftersom, hvorvidt}）与平行对称句式 \\
\midrule
2. 动词、\newline 时态语态 & 现在时（\textit{-r}）；情态动词（\textit{kan/skal/vil/må}）；规则与高频不规则过去时；祈使句；位置动词（\textit{ligger/står}） & \textbf{现在完成时（\textit{har/er} + 分词）}；\textbf{完成时 vs 过去时切换}；过去完成时；感官动词结构 & \textbf{被动语态（\textit{-s} 被动 vs \textit{blive} 被动）}；现在/过去分词定语；情态动词表主观推测 & \textbf{虚拟语气与非真实条件句}（\textit{havde... ville have...}）；动词名词化（\textit{-ning / -else}） & 学术/媒体引述动词立场光谱（\textit{påpege, hævde, fremhæve}）；趋势与统计动词体系 \\
\midrule
3. 名词、\newline 代词形容词 & 共性 \textit{en}/中性 \textit{et}；单复数定冠词后缀；物主代词；\textit{hvilken/hvilket/hvilke}；\textit{nogen/noget} & 形容词词尾（\textit{-t / -e}）与\textbf{比较级/最高级}；\textbf{反身物主代词 \textit{sin/sit/sine} vs \textit{hans/hendes}}；相互代词 \textit{hinanden} & 形容词名词化泛指（\textit{det vigtigste}）；复合名词构词与 \textit{fuge-s}；分数比例搭配 & 抽象名词与动词高阶固定搭配（\textit{stille krav om, tage stilling til}） & 丹麦文化与思辨核心抽象词辨析（\textit{dannelse vs uddannelse}） \\
\midrule
4. 连词与\newline 篇章衔接 & 并列连词（\textit{og, men, eller, for, så}）；顺序副词（\textit{først, bagefter, så}） & 核心从句连词（\textit{fordi, da vs når, hvis, selvom, mens}）；篇章副词（\textit{derfor, alligevel, desværre}） & 高级因果与让步（\textit{skyldes at, på trods af, mens derimod}）；比较结构（\textit{jo... desto...}） & 论证与转折衔接（\textit{ikke desto mindre, som følge heraf, medmindre, forudsat at}） & 思辨与辩证过渡（\textit{paradoksalt nok, i takt med at}） \\
\midrule
5. 语用、\newline 语音与文体 & 软 d、不发音辅音、吞音、重音与 \textit{stød} 感知；考场救场元语言（Metasprog） & 私人邮件格式、日常叙事节奏、多人物信息对比 & 客观说明文、社会制度介绍、采访转述语体 & \textbf{灵魂情态小品词（\textit{jo, da, nok, vist, lige, vel}）}；日常习语；正反辩论投稿 & \textbf{丹麦式反讽与低调委婉（\textit{Underdrivelse}）}；图表三步评论；多源综述 \\
\bottomrule
\end{tabularx}
}

% =====================================================================
\section{分册详细目录与每课语法靶点（共 104 课）}
% =====================================================================

\subsection{第一册：《新概念丹麦语 1 · Første Skridt》（共 24 课）}
\textbf{定位}：对标 NCE 1 + DU3 Modul 1（A1 $\rightarrow$ A2-），每课 60--100 词。前 5 单元对应 DU3 口试五大话题域，每单元前 3 课情景对话，第 4 课第一人称独白模板。

{\small
\renewcommand{\arraystretch}{1.28}
\begin{longtable}{@{} >{\raggedright\arraybackslash}p{2.65cm} >{\raggedright\arraybackslash}p{4.6cm} >{\raggedright\arraybackslash}p{5.15cm} >{\raggedright\arraybackslash}p{3.2cm} @{}}
\toprule
\textbf{单元} & \textbf{课号与标题（丹 / 中）} & \textbf{核心语法靶点（Key Structures）} & \textbf{词汇场景与语音重点} \\
\midrule
\endfirsthead
\toprule
\textbf{单元} & \textbf{课号与标题（丹 / 中）} & \textbf{核心语法靶点（Key Structures）} & \textbf{词汇场景与语音重点} \\
\midrule
\endhead
\bottomrule
\endfoot
\textbf{Unit 1}\newline 初到丹麦与\newline 语言学校\newline \textit{(Sprogskole)} & \textbf{L1} \textit{Undskyld, er det plads 14?}\newline （请问，这是 14 号座吗？） & 主格人称代词；系动词 \textit{at være}；\textbf{是非问句倒装（动词放句首）}；\textbf{\textit{ja / nej / jo} 应答区分}；简短回答（\textit{Det er det}） & 火车座位核对与寒暄；不发音的 d/t 与软 d 初探 \\
 & \textbf{L2} \textit{Hvem er du, og hvor kommer du fra?}\newline （你是谁，来自哪里？） & \textbf{\textit{hv-} 疑问句倒装}（\textit{hvad, hvor, hvem, hvordan}）；规则动词现在时（\textit{+r}） & 国籍、语言、个人信息；\textit{hv-} 词首 h 不发音 \\
 & \textbf{L3} \textit{På sprogskolen i første pause}\newline （语言学校的第一次课间休息） & \textbf{名词共性（n-ord \textit{en}）与中性（t-ord \textit{et}）}；单数定冠词后缀（\textit{-en / -et}）；\textit{Der er...} 存在句 & 语言学校设施（\textit{en kantine, et hold}）；定冠词后缀弱读 \\
 & \textbf{L4} \textit{Min sprogskole og mit hold} [独白]\newline （我的语言学校和班级） & 名词复数（\textit{-er/-e}）与复数定冠词（\textit{-erne/-ene}）；\textbf{\textit{hvor mange} vs \textit{hvor meget}}；\textbf{\textit{hvilken / hvilket / hvilke}} & 课程安排、学费、同学与老师；词尾弱读音 [ə] \\
\midrule
\textbf{Unit 2}\newline 日常作息与交通\newline \textit{(Tid \& Transport)} & \textbf{L5} \textit{Hvornår går bussen?}\newline （公交车几点开？） & 钟点与星期表达；时间疑问词（\textit{hvornår, hvor tit, hvor lang tid}）；\textbf{句首时间状语引发的主句 V2 倒装} & 公共交通与时刻表；数字与时间吞音连读 \\
 & \textbf{L6} \textit{En travl mandag morgen}\newline （忙碌的周一早晨） & \textbf{中心状语位置}（\textit{ikke, altid, aldrig, tit, også, kun} 在陈述句与倒装问句中的位置）；反身动词初探（\textit{skynder sig}） & 日常起居作息；高频副词轻重音规则 \\
 & \textbf{L7} \textit{Cykel eller metro til arbejde?}\newline （骑车还是坐地铁去上班？） & \textbf{情态动词}（\textit{kan, skal, vil, må} + 不带 \textit{at} 的动词原形）；交通介词搭配（\textit{med bus, på cykel, i bil}） & 通勤方式、停车与路线；情态动词弱读 \\
 & \textbf{L8} \textit{Fra morgen til aften} [独白]\newline （从早到晚的一天） & 顺序连接词（\textit{først, bagefter, så, til sidst}）引导的连续 V2 倒装句；日程短语（\textit{stå op, møde, få fri}） & 完整描述个人一日作息；句群节奏与停顿 \\
\midrule
\textbf{Unit 3}\newline 住房与邻里\newline \textit{(Bolig)} & \textbf{L9} \textit{En ny lejlighed til leje}\newline （一套出租的新公寓） & \textbf{形容词基础词尾变化}（n-ord 原形 / t-ord \textit{+t} / 复数与定指 \textit{+e}） & 租房广告、房间面积、楼层与租金 \\
 & \textbf{L10} \textit{Hvor står vaskemaskinen?}\newline （洗衣机放在哪儿？） & \textbf{四大静态位置动词}（\textit{ligger, står, hænger, sidder}）；动/静态方位副词（\textit{ind/inde, ud/ude, op/oppe}） & 家具电器、厨房与洗衣房；\textit{-ld/-nd/-rd} 不发音 d \\
 & \textbf{L11} \textit{Er der nogen hjemme?}\newline （有人在家吗？） & 不定代词 \textbf{\textit{nogen / noget / nogle}} 与否定表达 \textit{ikke nogen / ikke noget (ingen / intet)} & 拜访邻居、借东西、公寓设施问答 \\
 & \textbf{L12} \textit{Min bolig og mine naboer} [独白]\newline （我的住处与邻居） & 物主代词（\textit{min/mit/mine, vores}）及第三人称 \textbf{\textit{sin/sit/sine} vs \textit{hans/hendes}} 初探 & 描述居住环境、合租/独居与家务分工 \\
\midrule
\textbf{Unit 4}\newline 工作、实习与学习\newline \textit{(Arbejde \& Uddannelse)} & \textbf{L13} \textit{Første dag i bageriet}\newline （面包店上班第一天） & 祈使句（Imperativ）；泛指代词 \textbf{\textit{man}}；职责与规定（\textit{man skal, man må ikke, man behøver ikke}） & 工作服、排班、安全规定；喉塞音 \textit{stød} 感知 \\
 & \textbf{L14} \textit{Frokostpause i kantinen}\newline （食堂的午休时间） & 人称代词宾格（\textit{mig, dig, ham, hende, os, jer, dem}）；喜好表达（\textit{kan godt lide, er glad for}） & 同事闲聊、午餐文化、职场人际 \\
 & \textbf{L15} \textit{Et studie på universitetet}\newline （在大学读书） & \textbf{五大并列连词}（\textit{og, men, eller, for, så}）——重点区分 \textit{for}（后接正常语序）与 \textit{så}（“所以”后接倒装语序） & 专业课程、小组作业、考试与教学语言 \\
 & \textbf{L16} \textit{Mit arbejde og mine kolleger} [独白]\newline （我的工作与同事） & 职业名词零冠词规则（\textit{Jeg er ingeniør}）；时间频率介词（\textit{om dagen, om ugen, i timen}） & 完整介绍工作/实习/学业情况 \\
\midrule
\textbf{Unit 5}\newline 休闲、运动与消费\newline \textit{(Fritid \& Indkøb)} & \textbf{L17} \textit{Et medlemskab i fitnesscentret}\newline （办一张健身房会员卡） & 价格、折扣与报名句型（\textit{tilmelde sig, prøve gratis}）；不定式标志 \textit{at} 的使用与省略 & 俱乐部报名、开放时间、学生优惠 \\
 & \textbf{L18} \textit{Hvad lavede du i weekenden?}\newline （你周末做了什么？） & \textbf{规则动词一般过去时（Datid）}：第 1 类（\textit{-ede}）与第 2 类（\textit{-te}）；过去时间状语（\textit{i går, i søndags}） & 周末活动、聚会与家务；词尾 \textit{-ede} 发音 [əðə] \\
 & \textbf{L19} \textit{En tur i supermarkedet}\newline （去一趟超市） & \textbf{20 个最高频不规则动词过去时}（\textit{var, havde, gik, fik, tog, så, kom, købte, fandt}） & 超市采购、商品计量与结账对话 \\
 & \textbf{L20} \textit{Min fritidsaktivitet} [独白]\newline （我的业余爱好） & 运动与爱好固定搭配（\textit{gå til..., spille..., dyrke...}）；形式主语评价句（\textit{Det er sjovt/sundt at...}） & 完整描述一项个人业余活动 \\
\midrule
\textbf{Unit 6}\newline 日常交际与总复习\newline \textit{(Overblik)} & \textbf{L21} \textit{Hos lægen}\newline （看医生） & 身体部位与疼痛表达（\textit{Det gør ondt i...}）；建议与必要性（\textit{Du bør..., Det er vigtigt at...}） & 预约家庭医生、描述症状与请病假 \\
 & \textbf{L22} \textit{En invitation til fødselsdag}\newline （生日聚会邀请） & 丹麦语表达“将来”的三种方式（现在时 + 将来时间、\textit{skal} 表计划、\textit{vil} 表意愿） & 节日庆祝、接受与委婉拒绝邀请 \\
 & \textbf{L23} \textit{Undskyld, kan du gentage det?}\newline （抱歉，能重复一遍吗？） & \textbf{交际救场元语言（Metasprog）}；简短附和句（\textit{Det gør jeg også / Det gør jeg heller ikke}） & 听不懂时的追问、拼写确认与考场互动 \\
 & \textbf{L24} \textit{Mit første halve år i Danmark} [综合]\newline （我在丹麦的前半年） & \textbf{第一册语法总汇演}：过去与现在对比、V2 倒装与全套疑问句综合操练 & 个人跨文化生活小结（100 词） \\
\end{longtable}
}

\subsection{第二册：《新概念丹麦语 2 · Fra Sætning til Tekst》（共 24 课）}
\textbf{定位}：对标 NCE 2 前半 + DU3 Modul 2（A2），每课 100--180 词。全面对标 DU3 Modul 2 的邮件写作、口语信息差及阅读四大题型。

{\small
\renewcommand{\arraystretch}{1.28}
\begin{longtable}{@{} >{\raggedright\arraybackslash}p{2.65cm} >{\raggedright\arraybackslash}p{4.6cm} >{\raggedright\arraybackslash}p{5.15cm} >{\raggedright\arraybackslash}p{3.2cm} @{}}
\toprule
\textbf{单元} & \textbf{课号与标题（丹 / 中）} & \textbf{核心语法靶点（Key Structures）} & \textbf{对标体裁与 Modul 2 题型} \\
\midrule
\endfirsthead
\toprule
\textbf{单元} & \textbf{课号与标题（丹 / 中）} & \textbf{核心语法靶点（Key Structures）} & \textbf{对标体裁与 Modul 2 题型} \\
\midrule
\endhead
\bottomrule
\endfoot
\textbf{Unit 1}\newline 私人邮件与\newline 近期经历\newline \textit{(Skrivning \& Mundtlig)} & \textbf{L1} \textit{Et postkort fra Rom — og en e-mail hjem}\newline （来自罗马的明信片与回信） & \textbf{现在完成时（\textit{har/er} + 分词）}；助动词 \textit{har} vs \textit{er}；\textbf{完成时 vs 过去时黄金切换} & 幽默叙事 + 邮件写作：旅行见闻（对标 \textit{Skrivning Opg 1}） \\
 & \textbf{L2} \textit{En ny lille hund i familien}\newline （家里新来的小狗） & \textbf{原因从句 \textit{fordi} 与从句核心语序（主语 + \textit{ikke/altid/også} + 动词）} & 邮件写作：养宠物经历（对标 \textit{3\_2\_skrive\_A}） \\
 & \textbf{L3} \textit{En uforglemmelig koncert}\newline （一场难忘的音乐会） & \textbf{现在完成时 vs 一般过去时进阶实战}（完成时引出经历，过去时叙述具体细节） & 邮件写作：文化演出体验（对标 \textit{3\_2\_skrive\_B}） \\
 & \textbf{L4} \textit{Besøg fra udlandet}\newline （来自国外的访客） & 时间从句 \textbf{\textit{da}（过去一次性）vs \textit{når}（每次/将来）}；\textbf{“从句在前，主句必倒装”}铁律 & 邮件写作：接待亲友来访（对标 \textit{3\_2\_skrive\_C}） \\
 & \textbf{L5} \textit{Majas nye arbejde på caféen}\newline （玛雅在咖啡馆的新工作） & 第三人称信息问答与转述；\textbf{间接疑问从句}（\textit{Jeg ved ikke, om hun... / hvornår hun...}） & 口语信息差问答（对标 \textit{Opgave 2: Informationskløft}） \\
 & \textbf{L6} \textit{Vi er endelig flyttet!}\newline （我们终于搬家了！） & 时间跨度介词精析（\textbf{\textit{i to år} vs \textit{for to år siden} vs \textit{om to år}}）；书信开头与结尾套语 & 邮件写作：搬家与新生活近况 \\
\midrule
\textbf{Unit 2}\newline 健康生活与\newline 社区小故事\newline \textit{(Læsning Opg 1 \& 2)} & \textbf{L7} \textit{Vi cykler til arbejde i maj}\newline （五月我们骑车上班） & 目的不定式（\textbf{\textit{for at} + inf.}）；让步从句 \textbf{\textit{selvom}}（虽然……但是……） & 活动倡议与健康生活（对标 \textit{Læsning Opg 1}） \\
 & \textbf{L8} \textit{Kajak, karate eller gåklub?}\newline （皮划艇、空手道还是健步走？） & \textbf{形容词与副词的比较级与最高级}（\textit{-ere / -st} 及不规则变化 \textit{bedre, flere, mere}） & 运动俱乐部广告对比与推荐 \\
 & \textbf{L9} \textit{Peter Sørensen ringer til politiet}\newline （彼得·索伦森报警记） & 同时性时间从句 \textbf{\textit{mens}}；感官动词宾语补足结构（\textit{høre nogen larme}） & 悬念叙事与段落逻辑一致（对标 \textit{Læsning Opg 2}） \\
 & \textbf{L10} \textit{En sur seddel i vaskekælderen}\newline （洗衣房里的一张愤怒纸条） & \textbf{过去完成时（Førdatid：\textit{havde/var} + 分词）}；叙事中的“过去的过去”与因果回溯 & 邻里小摩擦与幽默生活故事 \\
 & \textbf{L11} \textit{Gode råd til en sundere hverdag}\newline （健康日常的好建议） & 条件从句 \textbf{\textit{hvis}}；成对连词 \textbf{\textit{både... og...}} 与 \textbf{\textit{hverken... eller...}} & 健康建议与生活习惯说明文 \\
 & \textbf{L12} \textit{Da toget stod stille i to timer}\newline （当火车停运两小时） & 形容词加 \textit{-t} 变方式副词 vs 程度副词规律；突发事件叙事弧线 & 日常意外经历叙事（NCE 2 经典风格） \\
\midrule
\textbf{Unit 3}\newline 居住选择与\newline 丹麦初体验\newline \textit{(Læsning Opg 3 \& 4)} & \textbf{L13} \textit{Tre studerende om deres bolig i Aarhus}\newline （三位大学生谈奥胡斯的住处） & \textbf{关系从句 \textit{som} 与 \textit{der}}（\textit{der} 仅作主语，\textit{som} 可作主语、宾语或搭配后置介词） & 多人物对比阅读（对标 \textit{Læsning Opg 4}） \\
 & \textbf{L14} \textit{Danskerne finder kærlighed på nettet}\newline （丹麦人在网上找真爱） & \textbf{高频篇章连接副词辨析}（\textit{derfor, alligevel, desværre, heldigvis, måske, desuden} 放句首倒装） & 社会现象与完形填空（对标 \textit{Læsning Opg 3}） \\
 & \textbf{L15} \textit{Det første år i Danmark}\newline （在丹麦的第一年） & 心理与适应动词搭配（\textit{vænne sig til, savne, glæde sig til, være nødt til}）；宾语从句 \textit{at...} & 三位外国人的丹麦适应记（对标 \textit{Læsning Opg 4}） \\
 & \textbf{L16} \textit{Fællesmiddag og rengøringsdag}\newline （公共晚餐与大扫除日） & 相互代词 \textbf{\textit{hinanden}}；共同与轮流表达（\textit{skiftes til at, være fælles om, hver vores}） & 丹麦合租与学生宿舍（Kollegium）文化 \\
 & \textbf{L17} \textit{Når familien vokser}\newline （当家庭迎来新成员） & \textbf{反身物主代词 \textit{sin/sit/sine} 进阶排雷}（复数主语、跨从句指代判定） & 家庭变迁、育儿假与跨国探亲 \\
 & \textbf{L18} \textit{En klage over larm i opgangen}\newline （一封关于楼道噪音的投诉信） & 半正式信函结构；委婉批评与请求句式（\textit{Jeg skriver til jer, fordi... / Vil I være søde at...}） & 给物业或房东写反馈邮件 \\
\midrule
\textbf{Unit 4}\newline 消费、职场与\newline 进阶叙事\newline \textit{(Modul 2 $\rightarrow$ 3)} & \textbf{L19} \textit{En pakke, der aldrig kom frem}\newline （一个从未送达的包裹） & 时间界限介词与连词（\textbf{\textit{siden, indtil, inden, før, i løbet af}}） & 网购维权与客服沟通故事 \\
 & \textbf{L20} \textit{Samtale med chefen om arbejdstid}\newline （与老板谈工作时间） & 礼貌虚拟：用过去时表委婉请求（\textit{Kunne vi måske...? / Jeg ville gerne høre, om...}） & 职场沟通、调班与休假申请 \\
 & \textbf{L21} \textit{Sommerhus ved Vesterhavet}\newline （北海边的度假小屋） & 丹麦语独特的\textbf{三维空间方位副词组合}（\textit{ud til kysten, op i huset, hjem hos os}） & 丹麦夏日度假屋文化与自然描写 \\
 & \textbf{L22} \textit{Hvordan overlever man den danske vinter?}\newline （如何熬过丹麦的冬天？） & \textbf{分裂强调句初探}（\textit{Det er mørket, der er sværest / Det er om vinteren, at...}） & 季节、日照与生活小窍门 \\
 & \textbf{L23} \textit{Et uventet møde på loppemarkedet}\newline （跳蚤市场上的意外重逢） & \textbf{直接引语转间接引语}（人称代词与时态后移呼应：现在时转过去时，完成时转过去完成时） & 趣味人物故事与对话转述 \\
 & \textbf{L24} \textit{Tilbageblik og fremtidsplaner} [综合]\newline （回顾与未来规划） & \textbf{第二册语法总汇演}：过去/完成/将来时态串联，主从复合句与篇章连接词综合写作 & 150--180 词长邮件与个人经历总结 \\
\end{longtable}
}

\subsection{第三册：《新概念丹麦语 3 · Hverdag og Samfund》（共 20 课）}
\textbf{定位}：对标 NCE 2 后半--NCE 3 前半 + DU3 Modul 3（B1），每课 200--300 词，聚焦丹麦社会制度、职场规则与文化心理。

{\small
\renewcommand{\arraystretch}{1.28}
\begin{longtable}{@{} >{\raggedright\arraybackslash}p{2.65cm} >{\raggedright\arraybackslash}p{5.7cm} >{\raggedright\arraybackslash}p{7.25cm} @{}}
\toprule
\textbf{单元} & \textbf{课号与标题（丹 / 中）} & \textbf{核心语法靶点（Key Structures）} \\
\midrule
\endfirsthead
\toprule
\textbf{单元} & \textbf{课号与标题（丹 / 中）} & \textbf{核心语法靶点（Key Structures）} \\
\midrule
\endhead
\bottomrule
\endfoot
\textbf{Unit 1}\newline 丹麦职场与\newline 教育体系 & \textbf{L1} \textit{Det flade hierarki på danske arbejdspladser}（丹麦职场的扁平化层级） & \textbf{被动语态核心区分：\textit{-s} 被动（客观规则/习惯/通用说法）vs \textit{blive} 被动（具体事件/动作过程）} \\
 & \textbf{L2} \textit{Fagforening, A-kasse og "flexicurity"}（工会、失业保险与“灵活安全”模式） & 介词 + 关系代词（\textit{hvilket, hvori, hvorfor}）与指代整句的关系从句（\textit{..., hvilket betyder, at...}） \\
 & \textbf{L3} \textit{Fra børnehave til universitet: Livslang læring}（从幼儿园到大学：终身学习） & \textbf{现在分词（\textit{-ende}）与过去分词（\textit{-et/-ede}）作定语形容词}（\textit{en spændende uddannelse, de uddannede unge}） \\
 & \textbf{L4} \textit{Når arbejdsdagen slutter klokken fire}（当下午四点下班铃响起） & \textbf{省略 \textit{hvis} 的条件倒装句}（\textit{Har man børn, skal man hente dem inden kl. 17}） \\
 & \textbf{L5} \textit{En jobsamtale i en dansk virksomhed}（一家丹麦企业的求职面试） & 职场核心动词 + 介词固定搭配（\textit{have ansvar for, tage initiativ til, sætte pris på, bidrage til}） \\
\midrule
\textbf{Unit 2}\newline 文化心理与\newline 社会默契 & \textbf{L6} \textit{Hvad er egentlig "hygge"?}（究竟什么是“Hygge”？） & 形容词名词化泛指结构（\textit{det hyggelige, det vigtigste, de unge, de ældre}） \\
 & \textbf{L7} \textit{Foreningsdanmark: Hvorfor er alle med i en klub?}（社团丹麦：为什么人人都参加俱乐部？） & 因果表达进阶（\textbf{\textit{skyldes at...} vs \textit{føre til at...} vs \textit{på grund af}}） \\
 & \textbf{L8} \textit{Janteloven i det moderne samfund}（现代社会中的《詹代法则》） & 情态动词的主观推测与传闻用法（\textit{må være, kan være, skulle efter sigende være}） \\
 & \textbf{L9} \textit{Fra vikingetid til velfærdsstat}（从维京时代到福利国家） & 历史叙事中的时态规范与长跨度时间衔接词（\textit{op gennem 1900-tallet, i efterkrigstiden}） \\
 & \textbf{L10} \textit{Tillid: Når barnevognen står ude på gaden}（信任：当婴儿车停在街边） & 高级让步与对比结构（\textbf{\textit{til trods for at..., på trods af} + 名词, \textit{mens... derimod}}） \\
\midrule
\textbf{Unit 3}\newline 家庭、消费与\newline 数字化 & \textbf{L11} \textit{MitID, e-Boks og det papirløse samfund}（MitID、电子邮箱与无纸化社会） & \textbf{丹麦语必考句法：介词 + 形式宾语 \textit{det} + \textit{at/hv-} 从句}（\textit{regne med, at... / være i tvivl om, hvorvidt...}） \\
 & \textbf{L12} \textit{Delingsøkonomi og genbrugstøj}（共享经济与二手服装） & 递进与等比比较句型（\textbf{\textit{jo} + 比较级..., \textit{desto/jo} + 比较级...}；\textit{lige så... som}） \\
 & \textbf{L13} \textit{Børneopdragelse og "curlingforældre"}（育儿观念与“冰壶父母”） & 客观引述与社会舆论表达（\textit{Ifølge eksperter..., Det diskuteres ofte, om..., Mange peger på, at...}） \\
 & \textbf{L14} \textit{Når kontanterne forsvinder}（当现金逐渐消失） & 趋势变化动词与程度副词（\textit{stige, falde, i stigende grad, stadigt flere/færre}） \\
 & \textbf{L15} \textit{Alene, sammen eller i regnbuefamilie?}（独居、同居还是多元家庭？） & 分数、比例与群体主谓搭配（\textit{hver tredje dansker, halvdelen af befolkningen, størstedelen af}） \\
\midrule
\textbf{Unit 4}\newline 环境、城市与\newline 公共医疗 & \textbf{L16} \textit{Vindmøller og drømmen om grøn energi}（风车与绿色能源之梦） & \textbf{复合名词构词法与连接字母 \textit{fuge-s / fuge-e}}：如何拆解与自造地道丹麦语长复合词 \\
 & \textbf{L17} \textit{Cykelstier i København: Mere end bare transport}（哥本哈根自行车道：不止于交通） & 分裂句进阶与排比结构（\textit{Ikke alene..., men også... / Det er først og fremmest..., der...}） \\
 & \textbf{L18} \textit{Det danske sundhedsvæsen under lup}（放大镜下的丹麦医疗体系） & 书面复合介词短语（\textbf{\textit{i forhold til, i forbindelse med, med hensyn til, på baggrund af}}） \\
 & \textbf{L19} \textit{Udkantsdanmark eller storbydrømme?}（边缘乡村还是大城市之梦？） & 城乡对比说明文的段落衔接架构（平行对比法 vs 逐点对比法） \\
 & \textbf{L20} \textit{At finde sin plads i et nyt sprog} [综合]（在新语言中找到自己的位置） & \textbf{第三册语法总汇演}：被动句、介词从句与 250--300 词社会话题说明文写作 \\
\end{longtable}
}

\subsection{第四册：《新概念丹麦语 4 · Holdning og Debat》（共 20 课）}
\textbf{定位}：对标 NCE 3 + DU3 Modul 4（B1+ $\rightarrow$ B2-），每课 300--450 词，聚焦情态小品词、报刊专栏与正反辩论。

{\small
\renewcommand{\arraystretch}{1.28}
\begin{longtable}{@{} >{\raggedright\arraybackslash}p{2.65cm} >{\raggedright\arraybackslash}p{5.7cm} >{\raggedright\arraybackslash}p{7.25cm} @{}}
\toprule
\textbf{单元} & \textbf{课号与标题（丹 / 中）} & \textbf{核心语法与修辞靶点（Key Structures \& Style）} \\
\midrule
\endfirsthead
\toprule
\textbf{单元} & \textbf{课号与标题（丹 / 中）} & \textbf{核心语法与修辞靶点（Key Structures \& Style）} \\
\midrule
\endhead
\bottomrule
\endfoot
\textbf{Unit 1}\newline 灵魂语气与\newline 语体转换 & \textbf{L1} \textit{Hvad mener danskerne egentlig, når de siger "det er fint"?}（丹麦人说“挺好”时到底是什么意思？） & \textbf{灵魂情态小品词 I：\textit{jo}（你我皆知）、\textit{da}（当然/略带反驳）、\textit{nok}（大概/恐怕）、\textit{vist}（听说/好像）}的语用潜台词与句中位置 \\
 & \textbf{L2} \textit{Skal vi lige tage en kop kaffe?}（我们要不喝杯咖啡？） & \textbf{灵魂情态小品词 II：\textit{lige}（稍微/顺便）、\textit{vel}（寻求认同“吧？”）、\textit{bare}、\textit{egentlig}、\textit{ellers}} 的语气软化功能 \\
 & \textbf{L3} \textit{Fra talesprog til skriftsprog: En nyhedsartikel bliver til}（从口语采访到书面报道） & 口语松散句（左移外置、口头禅）向\textbf{紧凑书面语（名词短语化、定语从句化）}的改写技巧 \\
 & \textbf{L4} \textit{Kunsten at være uenig på en høflig måde}（礼貌表达不同意见的艺术） & 让步与反驳高阶句式（\textit{Ganske vist..., men ikke desto mindre... / Det kan godt være, at..., men...}） \\
 & \textbf{L5} \textit{Faste udtryk og billedsprog i hverdagen}（日常习语与比喻表达） & 25 个高频丹麦语成语/习语在议论与叙事中的活用（如 \textit{have mange bolde i luften, slå to fluer med ét smæk}） \\
\midrule
\textbf{Unit 2}\newline 福利国家的\newline 困境与辩论 & \textbf{L6} \textit{Hvem skal passe de ældre i fremtiden?}（未来谁来照料老人？） & 政策预测与被动预期结构（\textit{forventes at, skønnes at, står over for en alvorlig udfordring}） \\
 & \textbf{L7} \textit{SU: En investering i fremtiden eller en dyr gavebod?}（助学金 SU：未来投资还是昂贵福利？） & \textbf{正反方观点对立框架}（\textit{Tilhængere fremhæver, at...} $\leftrightarrow$ \textit{Modstandere indvender derimod, at...}） \\
 & \textbf{L8} \textit{4-dages arbejdsuge: Produktivitet eller ønsketænkning?}（四天工作制：生产力还是空想？） & \textbf{非真实条件句与虚拟语气}（用过去完成时 + \textit{ville have} + 分词 表达与事实相反的假设推演） \\
 & \textbf{L9} \textit{Skat med glæde? Debatten om topskatten}（快乐纳税？关于最高边际税的争论） & 严谨因果链条副词（\textbf{\textit{som følge heraf, af den grund, følgelig, hvorved}}） \\
 & \textbf{L10} \textit{Integration og medborgerskab: Hvad kræver det?}（融入与公民身份：需要什么？） & 社会议题高频动宾搭配（\textit{stille krav om, tage stilling til, spille en afgørende rolle, bære præg af}） \\
\midrule
\textbf{Unit 3}\newline 科技、媒体与\newline 现代伦理 & \textbf{L11} \textit{Kunstig intelligens i klasseværelset: Hjælp eller snyd?}（课堂里的 AI：助手还是作弊？） & \textbf{动词名词化（Verbalsubstantiver med \textit{-ning / -else}）}：将冗长 \textit{når man...} 从句浓缩为正式名词短语 \\
 & \textbf{L12} \textit{Børn, skærmtid og sociale medier}（儿童、屏幕时间与社交媒体） & 限制与前提条件连词（\textbf{\textit{medmindre, forudsat at, såfremt, i det omfang}}） \\
 & \textbf{L13} \textit{Overvågning i det offentlige rum: Tryghed kontra privatliv}（公共监控：安全感 vs 隐私） & 价值权衡句型（\textit{afveje A over for B, gå på kompromis med, ske på bekostning af}） \\
 & \textbf{L14} \textit{Køb og smid væk-kulturen under pres}（受冲击的“买了就扔”消费文化） & 正式书面语中的\textbf{前置长分词定语}（如 \textit{den af regeringen vedtagne lov}）的识别与自然改写 \\
 & \textbf{L15} \textit{Nyheder i algoritmernes tidsalder: Ekkokamre og kildekritik}（算法时代的新闻：回音室与信源批判） & 批判性思维表达（\textit{sætte spørgsmålstegn ved, forholde sig kritisk til, nuancere debatten}） \\
\midrule
\textbf{Unit 4}\newline 教育、气候与\newline 全球视野 & \textbf{L16} \textit{Karakterræs og trivselskrisen blandt unge}（分数竞争与青少年心理危机） & 多重归因与相关性表达（\textit{hænge nøje sammen med, en kombination af flere faktorer}） \\
 & \textbf{L17} \textit{Klimamål og landbrugets grønne omstilling}（气候目标与农业绿色转型） & 政策举措与妥协词汇（\textit{iværksætte tiltag, stramme lovgivningen, indgå et bredt forlig}） \\
 & \textbf{L18} \textit{Turisme i København: Velsignelse eller byrde?}（哥本哈根旅游业：福音还是负担？） & \textbf{标准双刃剑议论文四段式模板}（背景引入 $\rightarrow$ 积极面 $\rightarrow$ 消极面 $\rightarrow$ 权衡结论与建议） \\
 & \textbf{L19} \textit{Dansk sprog i en engelsktalende verden}（英语世界中的丹麦语） & 设问推进与强调句（\textit{Spørgsmålet er imidlertid, hvorvidt... / Det står uden for enhver tvivl, at...}） \\
 & \textbf{L20} \textit{Kunsten at skrive et overbevisende debatindlæg} [综合]（如何写一篇有说服力的辩论投稿） & \textbf{第四册语法总汇演}：情态小品词、复杂从句与 350--400 词报刊读者来信（\textit{Debatindlæg}）写作 \\
\end{longtable}
}

\subsection{第五册：《新概念丹麦语 5 · Tanke og Perspektiv》（共 16 课）}
\textbf{定位}：对标 NCE 4 + DU3 Modul 5（PD3）及 Modul 6（Studieprøven，B2 $\rightarrow$ C1），每课 450--650 词，专攻图表评述、反讽修辞、伦理哲学思辨与学术综述。

{\small
\renewcommand{\arraystretch}{1.28}
\begin{longtable}{@{} >{\raggedright\arraybackslash}p{2.65cm} >{\raggedright\arraybackslash}p{5.7cm} >{\raggedright\arraybackslash}p{7.25cm} @{}}
\toprule
\textbf{单元} & \textbf{课号与标题（丹 / 中）} & \textbf{核心思辨、修辞与高阶语法靶点} \\
\midrule
\endfirsthead
\toprule
\textbf{单元} & \textbf{课号与标题（丹 / 中）} & \textbf{核心思辨、修辞与高阶语法靶点} \\
\midrule
\endhead
\bottomrule
\endfoot
\textbf{Unit 1}\newline 数据洞察与\newline 趋势分析\newline \textit{(直击 PD3 图表)} & \textbf{L1} \textit{Demografiens stille revolution: Når samfundet ældes}（人口结构的静默革命：当社会老龄化） & \textbf{统计图表全套高阶表达}：占比、峰值、拐点与差距变化（\textit{udgøre en andel på, toppe i, flade ud, en markant forskel}） \\
 & \textbf{L2} \textit{Deltidsarbejde og kønsroller på det danske arbejdsmarked}（丹麦劳动力市场上的兼职工作与性别角色） & \textbf{PD3 图表写作与口试黄金三步法}：\ding{172} 宏观概括数据趋势 $\rightarrow$ \ding{173} 剖析社会/制度成因 $\rightarrow$ \ding{174} 评估未来影响与对策 \\
 & \textbf{L3} \textit{Urbaniseringens pris: Hvem har råd til at bo i byen?}（城市化的代价：谁住得起大城市？） & 悖论与动态同步结构（\textbf{\textit{paradoksalt nok, i takt med at..., en selvforstærkende effekt}}） \\
 & \textbf{L4} \textit{Fritidsvaner i forandring: Fra foreningsliv til individuel træning}（变迁中的休闲：从社团生活到个人健身） & 社会学分析语汇：从具体生活现象提炼宏观社会趋势（\textit{individualisering, sammenhængskraft}） \\
\midrule
\textbf{Unit 2}\newline 修辞、反讽与\newline 媒体评论 & \textbf{L5} \textit{Dansk ironi og kunsten at underdrive}（丹麦式反讽与低调委婉的艺术） & \textbf{丹麦语修辞精髓：低调陈述/反语（Underdrivelse/Litote）}（如 \textit{ikke helt dårligt, en smule udfordrende}）的语境识别与运用 \\
 & \textbf{L6} \textit{Lederen i søndagsavisen: Velfærd eller skattelettelser?}（周日版社论：福利优先还是减税优先？） & 报刊社论（\textit{Leder}）的高阶让步句式（\textit{Hvorvidt..., står hen i det uvisse / Uanset hvordan man vender og drejer det...}） \\
 & \textbf{L7} \textit{Ekspertvælde eller sund fornuft?}（专家治国还是常识至上？） & \textbf{引述动词的态度立场光谱}：客观中立（\textit{konstatere, påpege}）vs 保持距离/质疑（\textit{hævde, påstå}）vs 认同强调（\textit{understrege, slå fast}） \\
 & \textbf{L8} \textit{Kulturkamp ved middagsbordet: Mad, moral og klimaskam}（餐桌上的文化之争：饮食、道德与气候羞耻） & 三层以上主从复合句的语序拆解、插入语标点规范与长句节奏控制 \\
\midrule
\textbf{Unit 3}\newline 伦理、制度与\newline 现代性思辨 & \textbf{L9} \textit{Ytringsfrihedens grænser i en polariseret tid}（极化时代的言论自由边界） & 原则与实践的辩证论证结构（\textit{principielt set} $\leftrightarrow$ \textit{i praksis, trække en hårfager grænse mellem A og B}） \\
 & \textbf{L10} \textit{Præstationssamfundet og jagten på det perfekte liv}（绩效社会与对完美生活的追逐） & 哲学与心理学思辨评论：用连续设问与概念拆解推动深度论述 \\
 & \textbf{L11} \textit{Folkeskolen mellem faglighed og dannelse}（公立学校在“学科能力”与“全人教养”之间） & \textbf{丹麦文化核心不可译词深度辨析}：\textit{dannelse} vs \textit{uddannelse}、\textit{tryghed} vs \textit{sikkerhed}、\textit{fællesskab} vs \textit{samfund} \\
 & \textbf{L12} \textit{Bioetik og fremtidens medicin: Hvor langt bør vi gå?}（生命伦理与未来医学：我们应当走多远？） & 伦理两难推演与警示句型（\textit{rejse fundamentale etiske dilemmaer, skabe præcedens for, uoverskuelige konsekvenser}） \\
\midrule
\textbf{Unit 4}\newline 文学、历史与\newline 高阶学术\newline \textit{(衔接 C1)} & \textbf{L13} \textit{H.C. Andersens skygge: Eventyr for børn eller satire for voksne?}（安徒生的影子：儿童童话还是成人讽刺？） & 文学批评与文化解读语体：双关、隐喻与文本潜台词分析 \\
 & \textbf{L14} \textit{Danmark i Europa: Småstatsmentalitet og globalt udsyn}（欧洲中的丹麦：小国心态与全球视野） & 高阶书面连词与代词（\textbf{\textit{eftersom, idet, hvorvidt, desuagtet, herunder}}）在历史政论文中的运用 \\
 & \textbf{L15} \textit{Det gode liv: Lykkeforskning og det skandinaviske paradoks}（美好生活：幸福感研究与斯堪的纳维亚悖论） & \textbf{多源观点综合论述（Syntese）}：将两篇不同立场的背景材料整合为一篇批判性学术长文 \\
 & \textbf{L16} \textit{Sproget som hjem: Tanker om at leve på to sprog} [终章]（以语言为家：关于双语生活的哲思） & \textbf{全五册总收官}：兼顾文学美感与思辨深度的散文体议论文 \\
\end{longtable}
}

% =====================================================================
\section{单课标准编写规范（全书统一 6 大模块）}
% =====================================================================

为确保每课可独立完成“输入—精讲—操练—输出”的自学闭环，每一课统一包含以下 6 个板块：

\begin{enumerate}
  \item \textbf{学习导引（Oversigt \& Studietips）}：明确本课核心语法靶点、对标的 DU3 考点与建议用时。
  \item \textbf{一、Før du læser（导读提问）}：在读课文前先抛出 1--2 个问题，训练带着目的精读抓重点。
  \item \textbf{二、Tekst \& Udtale（课文、发音标注与译文）}：第 1--2 册配备 \textbf{IPA 国际音标（Charis SIL）}、软 d、不发音辅音、喉塞音 [ˀ] 及语流吞音提示，并附地道中文对照。
  \item \textbf{三、Gloser（核心生词表）}：\textbf{强制标注完整屈折形态}，统一格式为：
    \begin{itemize}
      \item 名词：\textit{en lejlighed (-en, -er, -erne)}
      \item 动词：\textit{at finde (finder, fandt, har fundet)}（完成时用 \textit{er} 的动词显式标出 \texttt{[er]}）
      \item 形容词：\textit{billig (billigt, billige)}
    \end{itemize}
  \item \textbf{四、Grammatik \& Faste vendinger（核心语法与地道表达精讲）}：每课讲透 1--3 个关联最紧密的语法规则 + 高频固定搭配。
  \item \textbf{五、Mønsterøvelser（句型操练） \& 六、Resumé / Skrivning（复述/写作实战与参考答案）}：包含语序转换、选词填空、中译丹，以及直接对标 DU3 口试或写作的微输出任务（附完整答案与解析）。
\end{enumerate}

\end{document}
